\documentclass[11pt,a4paper]{article}
\usepackage[T1]{fontenc}
\usepackage{lmodern,amsmath,amssymb}
\usepackage[margin=27mm]{geometry}
\usepackage[hidelinks]{hyperref}
\hypersetup{pdftitle={Width bounds, actual algebraic interfaces, and extremal lex growth},pdfauthor={}}
\setlength{\parskip}{3pt}
\setlength{\parindent}{1em}
\setlength{\emergencystretch}{2em}
\newtheorem{theorem}{Theorem}[section]
\newtheorem{lemma}[theorem]{Lemma}
\newtheorem{proposition}[theorem]{Proposition}
\newtheorem{remark}[theorem]{Remark}
\newenvironment{proof}[1][Proof]{\par\noindent\textit{#1.}\ \ignorespaces}{\unskip\nobreak\hfill\ensuremath{\square}\par}
\newcommand{\NN}{\mathbb N}
\newcommand{\kk}{\Bbbk}
\newcommand{\HF}{\operatorname{HF}}
\newcommand{\HS}{\operatorname{HS}}
\newcommand{\Lex}{\operatorname{Lex}}
\newcommand{\len}{\operatorname{length}}
\newcommand{\llbracket}{[\![}
\newcommand{\rrbracket}{]\!]}
\title{Width bounds, actual algebraic interfaces, and extremal lex growth}
\author{}
\date{Research delivery v0.3 --- 29 September 2026}
\begin{document}
\maketitle
\vspace{-1.4em}
\begin{center}\small
For independent mathematical review. Novelty has not been established.
Formal scope and remaining algebraic gaps are stated in Section~\ref{sec:formal}.
\end{center}
\begin{abstract}
This report closes a research-delivery phase on width bounds for numerical
semigroup rings with four minimal generators. A terminal column-square estimate
proves strict first and non-strict higher binomial bounds in the lex model for
all widths $w\ge4$. Published algebraic reductions give the semigroup application
as a conventional proof. Lean now also verifies actual monomial quotient spaces,
the least number of arbitrary polynomial generators, the genuine fiber $I/\mathfrak mI$,
and its identification with standard first Tor in the specified argument order.
An explicit family of finite-colength lex ideals realizes the matching growth
order: the attained maximum of the two-variable quotient dimension in the
budget class is $\Theta(w\log w)$. The higher Tor formulas and the complete
semigroup reduction are not formalized end to end. Novelty and human peer review
remain unestablished. Earlier drafts are preserved unchanged.
\end{abstract}

\section{The semigroup statement: conventional proof}
A numerical semigroup is an additive submonoid $\Gamma\subseteq\NN$ with finite complement, where $\NN=\{0,1,\ldots\}$. Suppose $\Gamma$ has exactly four minimal generators $g_0<g_1<g_2<g_3$, and put $m=g_0$ and $w=g_3-g_0$. For a field $\kk$, write
\[
 R_\Gamma=\kk\llbracket t^{g_0},t^{g_1},t^{g_2},t^{g_3}\rrbracket.
\]
All Betti numbers of $R_\Gamma$ are over its minimal regular presentation $\kk\llbracket X_0,X_1,X_2,X_3\rrbracket$.

\begin{theorem}\label{thm:main}
For every such $\Gamma$ and every field $\kk$,
\[
 \beta_i(R_\Gamma)\le i\binom{w+1}{i+1}\qquad(i\ge1).
\]
If $w\ge4$, then in addition
\[
 \beta_1(R_\Gamma)\le\binom{w+1}{2}-1.
\]
\end{theorem}

Caviglia--Moscariello--Sammartano (CMS) prove the non-strict bounds for $w\ge40$ in \cite[Theorem~1.5]{CMS}; their Remark~5.2 discusses the remaining small-width lex problem. Draft v0.1 treated that interval using a finite envelope certificate. The analytic proof works uniformly for all $w\ge4$ and makes the first bound strict. The case $w=3$ remains the known consecutive-generator case. We do not claim these stronger conclusions for arbitrary tangent cones or for arbitrary embedding dimension.

\section{A terminal budget for lex columns}
Let $S=\kk[x,y,z]$, ordered lexicographically with $x>y>z$, and let $L\subseteq(x,y,z)^2$ be an artinian lex ideal. Its standard monomials are those outside $L$. Assume that, for an integer $w\ge4$,
\begin{equation}\label{eq:budget}
 \sum_{t=0}^{d}\HF(S/L,t)\le1+dw\qquad(d\in\NN).
\end{equation}
Put $K=L\cap\kk[x,y]$, and let $\alpha$ be the least exponent such that $x^\alpha\in L$. Then $\alpha\ge2$. For $0\le i<\alpha$, define the positive column length
\[
 n_i=\#\{j\ge0:x^iy^j\notin K\},\qquad
 \ell=\len(\kk[x,y]/K)=\sum_{i=0}^{\alpha-1}n_i.
\]
Each column is an initial interval in its $y$ exponent. Since $x^iy^{n_i-1}$ is standard, lex order makes $y^{i+n_i-1}$ standard. Hence
\begin{equation}\label{eq:column-end}
 i+n_i\le n_0\qquad(0\le i<\alpha).
\end{equation}

\begin{lemma}\label{lem:terminal}
The column lengths satisfy
\begin{equation}\label{eq:terminal}
 \sum_{i=0}^{\alpha-1}\frac{n_i(n_i+1)}2\le1+(n_0-1)w.
\end{equation}
\end{lemma}
\begin{proof}
For fixed $i$, all monomials $x^iy^jz^k$ with $j+k<n_i$ are standard: first $x^iy^{j+k}$ is standard, and replacing some $y$ factors by $z$ lowers lex order. They have total degree at most $i+n_i-1\le n_0-1$, by \eqref{eq:column-end}. Their number is $n_i(n_i+1)/2$, and the sets for distinct $i$ are disjoint. Apply \eqref{eq:budget} once, at $d=n_0-1$.
\end{proof}

\begin{lemma}\label{lem:alpha}
One has $\alpha\le w-2$.
\end{lemma}
\begin{proof}
All monomials of degree below $\alpha$ are standard, since the largest such monomials $x^d$ are standard. Thus
\[
 \binom{\alpha+2}{3}\le1+(\alpha-1)w.
\]
Using
\[
 \alpha(\alpha+1)(\alpha+2)-6
   =(\alpha-1)(\alpha^2+4\alpha+6)
\]
and $\alpha\ge2$ gives $\alpha^2+4\alpha+6\le6w$. If $\alpha\ge w-1$, this implies $w^2+2w+3\le6w$, or $(w-1)(w-3)\le0$, impossible for $w\ge4$. The integer parameters therefore satisfy $\alpha\le w-2$.
\end{proof}

\section{A uniform strict estimate}
\begin{lemma}\label{lem:cauchy}
Under \eqref{eq:terminal},
\begin{equation}\label{eq:cauchy}
 (2\ell+\alpha-2w)^2\le
 \alpha\bigl[4(w-1)(w-2)+\alpha\bigr].
\end{equation}
\end{lemma}
\begin{proof}
Use the signed integer vector
\[
 v_0=2n_0+1-2w,\qquad v_i=2n_i+1\quad(1\le i<\alpha).
\]
There is no nonnegativity assumption on $v_0$. Direct expansion and \eqref{eq:terminal} yield
\begin{align*}
 \sum_i v_i&=2\ell+\alpha-2w,\\
 \sum_i v_i^2
 &=4\sum_i n_i(n_i+1)+\alpha-8wn_0-4w+4w^2\\
 &\le4(w-1)(w-2)+\alpha.
\end{align*}
Now apply $(\sum_i v_i)^2\le\alpha\sum_i v_i^2$. Equivalently, this step follows from the integer identity
\[
 \alpha\sum_i v_i^2-(\sum_i v_i)^2=\sum_{i<j}(v_i-v_j)^2\ge0.
\]
\end{proof}

\begin{theorem}\label{thm:strict}
For every integer $w\ge4$ under the hypotheses of Section~2,
\begin{equation}\label{eq:strict}
 \alpha+1+\ell\le\binom{w+1}{2}-1.
\end{equation}
Moreover,
\begin{equation}\label{eq:other}
 \alpha+2\ell\le2\binom{w+1}{3},\qquad
 \ell\le3\binom{w+1}{4}.
\end{equation}
\end{theorem}
\begin{proof}
Set $T=w^2-w-\alpha-2$. Lemma~\ref{lem:alpha} gives $T\ge w(w-2)>0$. A polynomial identity and the same lemma show
\begin{align*}
 T^2-\alpha\bigl[4(w-1)(w-2)+\alpha\bigr]
 &=(w-2)\bigl[(w-2)(w+1)^2-2\alpha(3w-1)\bigr]\\
 &\ge(w-2)^2(w-1)(w-3)>0.
\end{align*}
Combining this with \eqref{eq:cauchy} gives $2\ell+\alpha-2w<T$, since $T>0$. Rearranging yields
\[
 \alpha+1+\ell<\frac{w(w+1)}2,
\]
which is \eqref{eq:strict} by integrality.

Write $C=\binom{w+1}{2}$. Then $\ell\le C-\alpha-2$, so $\alpha+2\ell\le2C-\alpha-4\le2C$ and $\ell\le C$. Finally,
\[
 C\le\binom{w+1}{3},\qquad C\le3\binom{w+1}{4}\qquad(w\ge4),
\]
as is immediate from the ratios $(w-1)/3$ and $(w-1)(w-2)/4$. This proves \eqref{eq:other} without a bounded parameter search.
\end{proof}

\begin{remark}
The retained estimate \eqref{eq:cauchy} also gives
\[
 \ell\le w-\frac\alpha2+
   \sqrt{\alpha(w-1)(w-2)+\frac{\alpha^2}{4}}.
\]
Since $\alpha^2+4\alpha+6\le6w$, this also yields a conventional $O(w^{5/4})$ estimate. Section~\ref{sec:log} gives the stronger, formally verified logarithmic estimate; the square-root display is retained to explain the intermediate route.
\end{remark}

\section{The algebraic transfer: conventional layer}
For the finite-colength lex ideal $L$, the two-variable lex ideal $K$ has $\alpha+1$ minimal generators and rank one, so its Betti numbers as an ideal are $(\alpha+1,\alpha)$. The finite-colength formula \cite[equation~(7)]{CMS} gives
\begin{equation}\label{eq:betti}
 \beta_0^S(L)=\alpha+1+\ell,\quad
 \beta_1^S(L)=\alpha+2\ell,\quad
 \beta_2^S(L)=\ell.
\end{equation}
Thus Theorem~\ref{thm:strict} bounds the three positive Betti numbers of $S/L$, with the usual ideal-to-quotient shift by one.

\begin{proof}[Proof of Theorem~\ref{thm:main}]
The four-generator hypothesis gives $w\ge3$. If $w=3$, the generators are consecutive. Proposition~2.7 of \cite{HS}, with $r=4$, gives the desired bound for the graded semigroup ring. Localizing its minimal weighted-graded free resolution at the homogeneous maximal ideal and then completing preserve exactness and minimality, hence preserve the Betti numbers.

Suppose now $w\ge4$. CMS's construction starts with the Artinian reduction $R_\Gamma/(t^m)$ and passes to its associated graded ring and then to a homogeneous monomial initial ideal $J\subseteq Q=\kk[x,y,z]$. It has colength $m$ and no linear forms. The regular parameter $t^m$, the initial-ideal comparison and \cite[equation~(1)]{CMS} give
\[
 \beta_i(R_\Gamma)\le\beta_i^Q(Q/J)\quad(i\ge1).
\]
If $w\le m-2$, Theorem~3.3 of \cite{CMS} gives $\HS(Q/J,d)\le1+dw$. If $w\ge m-1$, the same inequality follows directly: the degree-zero contribution is one, and for $d\ge1$,
\[
 \HS(Q/J,d)\le m\le w+1\le1+dw.
\]
Thus the budget holds for every $w\ge4$, without invoking a separate large-width theorem.

Let $L=\Lex(J)$. It remains artinian, has no linear forms and has the same Hilbert function. The characteristic-independent Bigatti--Hulett--Pardue comparison \cite[Theorem~2.1]{CMS} implies
\[
 \beta_i(R_\Gamma)\le\beta_i^Q(Q/J)\le\beta_i^Q(Q/L).
\]
Apply \eqref{eq:betti} and Theorem~\ref{thm:strict} for $i=1,2,3$. For $i\ge4$, the Betti numbers vanish: $R_\Gamma$ is one-dimensional Cohen--Macaulay with a four-dimensional minimal regular presentation, and therefore has projective dimension three. This proves both conclusions.
\end{proof}

The Artinian assumption in \eqref{eq:betti} is not inferred from the combinatorial theorem. It comes from the semigroup reduction. For example, $(x^2,xy,xz,y^2,yz)$ retains infinitely many standard powers of $z$ and does not satisfy the finite-colength Betti formula. The general lex statement at $w=3$ is also false; that boundary is handled separately for semigroups.


\section{Actual ideals, generator numbers, and standard Tor}\label{sec:actual}
The following results describe the present machine-checked algebraic layer.
Write $A=\kk[x,y,z]$ and $\mathfrak m=(x,y,z)$, where $\kk$ is any field.
A monomial ideal is represented by an exponent upper set $E\subseteq\NN^3$;
membership of a monomial in the ideal is proved equivalent to membership of its
exponent in $E$. Lex closure refers to each fixed total degree, with $x>y>z$.
The notation $\ell$ now also denotes the dimension of the actual image of
$\kk[x,y]$ in $A/I$, rather than a count awaiting interpretation.

\begin{theorem}[Actual quotient dimensions and generator number]\label{thm:actual}
Let $I$ be a finite-colength lex monomial ideal, with $x\notin I$, satisfying
\eqref{eq:budget} for $w\ge4$. Let $\alpha$ be the least pure-$x$ exponent in $I$.
Then the standard monomial classes form a basis of $A/I$; the dimension of each
homogeneous image is its standard-monomial count. The full $xy$ image has
$\kk$-dimension $\ell=\sum_{i<\alpha}n_i$.
There is an explicit complete set of divisibility-minimal monomial generators
of $I$, of cardinality $\alpha+1+\ell$. For every finite set $P\subset A$ with
$(P)=I$, including nonhomogeneous polynomial sets,
\[
 |P|\ge\alpha+1+\ell.
\]
Equality is attained by that monomial set. Thus the actual least generator
number satisfies $\mu_A(I)=\alpha+1+\ell$, and obeys \eqref{eq:strict}.
\end{theorem}

The quotient bases come from explicit coefficient maps with the actual
monomial ideal as kernel. Finiteness of the quotient is equivalent to finiteness
of the set of standard exponents; for lex ideals it is also equivalent to the
presence of a pure power of $z$. The lower bound for arbitrary generating sets
uses the coefficients at all minimal exponents. Multiplication on those
coordinates acts through the constant coefficient. A generating family
therefore spans the whole coordinate space, forcing the cardinality bound.
This proves a least cardinality, rather than merely deletion-irredundance.

\begin{theorem}[Genuine fiber and standard first Tor]\label{thm:tor}
For any ideal $I\subseteq A$ there is a $\kk$-linear isomorphism
\[
 (A/\mathfrak m)\otimes_A I\simeq I/\mathfrak mI,
\]
and $A/\mathfrak m\simeq\kk$ as $\kk$-algebras. If $I\subseteq\mathfrak m$, then
\[
 \operatorname{Tor}_1^A(A/\mathfrak m,A/I)
 \simeq (A/\mathfrak m)\otimes_A I\simeq I/\mathfrak mI.
\]
For $I$ as in Theorem~\ref{thm:actual}, these spaces have dimension
$\alpha+1+\ell=\mu_A(I)$. Their bases correspond to the minimal monomial
classes and the pure tensors $1\otimes g$.
\end{theorem}

The first Tor here is mathlib's standard categorical Tor: the first argument
$A/\mathfrak m$ is fixed and the functor is derived in the second argument.
No interchange of the two arguments is being assumed as a formal theorem.
The construction uses the actual short exact sequence
$0\to I\to A\to A/I\to0$, a projective resolution starting with the specified
surjection, and the standard left-derived comparison. It does not assume that
$I$ is projective. Restriction to $\kk$ uses the original algebra map and its
scalar action, not an action defined by transport across the desired isomorphism.

The condition $I\subseteq\mathfrak m$ in the Tor statement is essential:
$I=A$ has nonzero fiber $A/\mathfrak m$ but
$\operatorname{Tor}_1^A(A/\mathfrak m,0)=0$.
Likewise, the equality between fiber dimension and the global minimum number
of generators is asserted here with the complete minimal monomial hypotheses;
it is not a theorem about arbitrary polynomial ideals.

\section{Harmonic growth and a realized lower family}\label{sec:log}
Let $H_n=\sum_{r=1}^{n}1/r$, with $H_0=0$. All inequalities involving
$H_n$ were first proved over the rationals, with exact natural-number counts.
For every ideal under Theorem~\ref{thm:actual}, the formal development gives
\begin{equation}\label{eq:harmonic}
 \ell\le3w+(2w-1)H_{2w-1},\qquad
 \mu_A(I)\le4w-1+(2w-1)H_{2w-1}.
\end{equation}
The second inequality uses $\alpha\le w-2$; it also bounds the dimension of the
standard Tor space in Theorem~\ref{thm:tor}. Comparison with the library harmonic
numbers and the natural logarithm yields
\begin{equation}\label{eq:logupper}
 \ell\le 3w+(2w-1)(1+\log(2w-1))\le10w\log w
 \qquad(w\ge4).
\end{equation}
The constant 10 is explicit and convenient; no optimality is claimed.

For $s\ge2$, define the following actual monomial ideal:
\[
 I_s=(x^iy^jz^k:(i+1)(i+j+k+1)>s^2).
\]
The exponent condition is upward closed and lex in each degree. The quotient
is finite dimensional: all monomials of degree at least $s^2$ belong to $I_s$.
The least pure-$x$ exponent is $s$. Its degree-$d$ $xy$ count is
\[
 h_d=\min\left(d+1,\left\lfloor\frac{s^2}{d+1}\right\rfloor\right).
\]
For the three-variable degree count $H^{(s)}_d$, the injection
$(i,j)\mapsto i(d+1)+j$ on standard exponent triples $(i,j,d-i-j)$
proves $H^{(s)}_d\le s^2$; also $H^{(s)}_0=1$. Consequently
\[
 \sum_{t=0}^{d}H^{(s)}_t\le1+ds^2\le1+dw
 \quad\text{whenever }s^2\le w.
\]
These are dimensions of the actual quotient homogeneous images by
Theorem~\ref{thm:actual}'s interface, not a proposed Hilbert sequence without
an ideal realizing it.

\begin{theorem}[Exact lower-family invariants]\label{thm:lower}
The actual $xy$ image for $I_s$ has dimension
\[
 \ell_s=\sum_{i=0}^{s-1}
 \left(\left\lfloor\frac{s^2}{i+1}\right\rfloor-i\right),\qquad
 s^2H_s\le\ell_s+\frac{s^2+s}{2}.
\]
Its least number of polynomial generators and its standard first Tor dimension
are both $s+1+\ell_s$.
\end{theorem}
Each summand is nonnegative; the implementation accounts for truncated natural
subtraction explicitly. This family supplies a lower bound in the abstract
lex budget class. No realization of these ideals by numerical semigroups is
claimed.

\section{An attained extremum and standard Theta}\label{sec:theta}
For each natural number $w$, let $\mathcal C_{\kk,w}$ consist of actual monomial
ideals $I$ of $A=\kk[x,y,z]$ such that:
\begin{enumerate}
\item the defining exponent set is upward closed and lex;
\item $A/I$ is finite dimensional over $\kk$ and $x\notin I$;
\item every degree satisfies the cumulative actual-dimension budget
$\sum_{t=0}^{d}\dim_{\kk}(A/I)_t\le1+dw$.
\end{enumerate}
Define $M_{\kk}(w)$ as the natural supremum of the dimensions of the actual
$xy$ images for ideals in this class. For $w\ge4$, $I_2$ makes the class nonempty,
and the strict binomial estimate bounds its set of dimensions. A nonempty
bounded set of natural numbers has a greatest member. Thus this supremum is
attained by an actual ideal, rather than an independent sum of optimal
per-degree envelopes.

\begin{theorem}[Formal extremal growth]\label{thm:theta}
For any field $\kk$, the maximum $M_{\kk}(w)$ is attained for $w\ge4$, and
\[
 \frac{w\log w}{16}\le M_{\kk}(w)\le10w\log w\qquad(w\ge64).
\]
As a function on natural widths with values in $\mathbb R$,
$M_{\kk}(w)=\Theta(w\log w)$ at infinity, in the standard mathlib
\texttt{Asymptotics.IsTheta} sense.
\end{theorem}
\begin{proof}[Proof structure]
The upper estimate is \eqref{eq:logupper} applied to a maximizing ideal.
For each $w\ge64$, choose $s=\lfloor\sqrt w\rfloor$. Then $s\ge2$ and
$s^2\le w$, so $I_s\in\mathcal C_{\kk,w}$. The exact floor-sum estimate in
Theorem~\ref{thm:lower}, harmonic-log comparison, and integer square-root
bounds give $\ell_s\ge w\log w/16$. All rounding and scalar conversions
are checked in Lean. This proves the lower estimate for every sufficiently
large integer $w$, not just squares. The two inequalities give the two standard
Big-O relations, including the required norms and eventual threshold.
\end{proof}
The real logarithm is used only for numerical bounds on dimensions. The
coefficient field need not have characteristic zero. The theorem describes
an extremum over the stated class; it is not a growth lower bound for each
member, nor for numerical semigroup Betti numbers. The upper comparison from
semigroups to lex ideals cannot be reversed to obtain such a lower bound.

\section{Formal verification, evidence and limits}\label{sec:formal}
This version records the final research boundary rather than claiming a full
formalization of Theorem~\ref{thm:main}. The conventional proof of that theorem
in Section~4 remains an internally audited algebraic argument using the cited
literature. Higher Tor/Betti formulas, the complete semigroup reduction,
the $w=3$ transfer and the final high-degree vanishing have not been connected
end to end in Lean. In \eqref{eq:betti}, the first generator count now has the
actual formal interfaces of Theorems~\ref{thm:actual}--\ref{thm:tor}; the higher
homological identities remain in the conventional layer.

\begin{center}
\small
\begin{tabular}{p{0.45\linewidth}p{0.45\linewidth}}
\hline
Result & Status at delivery\\ \hline
Uniform strict combinatorial bounds & Lean checked; no finite-certificate dependency\\
Actual monomial quotient dimensions & Lean checked\\
Least arbitrary polynomial generator number & Lean checked, stated monomial class\\
Actual fiber, residue field, standard Tor$_1$ & Lean checked, with $I\subseteq\mathfrak m$ for Tor\\
Explicit lower family and exact floor sum & Lean checked\\
Attained maximum and standard $\Theta(w\log w)$ & Lean checked, abstract lex budget class\\
Semigroup width theorem and higher Betti transfer & Conventional proof, internally reviewed\\
Finite examples (267 parameter pairs; 1807 semigroups) & Diagnostic experiments, not a general proof\\
Novelty and external expert review & Unestablished; external review deferred\\ \hline
\end{tabular}
\end{center}

\paragraph{Source entry points.}
{\raggedright
The project entry is \path{lean/WidthBounds.lean}. The main modules are
\texttt{AllWidths}, \texttt{TorOneBounds}, \texttt{LowerConstructionBounds},
and \texttt{LexGrowthTheta}; all are under \path{lean/WidthBounds/}.
Exact declarations, supporting modules, hypotheses and evidence are mapped in
\path{research/PROJECT_REPORT.md}. The standard asymptotic endpoint is
\texttt{maxLexLength\_isTheta} in namespace
\texttt{WidthBounds.LexExtremal}. The declaration source is authoritative.\par}

\paragraph{Verification record.}
Lean and mathlib are fixed at 4.22.0. The latest successful unified build is
\path{results/lean_lex_theta_build.txt}. Its SHA-256 is
\begin{center}\footnotesize\ttfamily
0b43de4c75bf3eeedc6af52ca63417c53\par
aae798e6593aefda2255eea40d230b8.
\end{center}
The two displayed lines concatenate to one digest (the final full stop is not
part of it). Accepted source/log hashes are recorded in
\path{research/claims.json}. This delivery changes documentation and
packaging only; it rechecks correspondence with the accepted sources instead
of claiming a fresh Lean build. To reproduce the formal checks, run
\texttt{lake build} in \texttt{lean/} with the pinned dependency manifest.
The delivery guide gives environment and artifact checks separately.

{\raggedright
Axiom output for the final twelve declarations contains only
\texttt{propext}, \texttt{Classical.choice}, and \texttt{Quot.sound}.
The project disallows proof placeholders, custom axioms and
\texttt{native\_decide}. Actual declaration-dependency audits of the new roots
exclude the old finite/small-width certificates and \texttt{sorryAx}.
These facts describe the checked declarations and their dependencies, not a
human endorsement of the intended mathematical model. Independent internal
AI reviews inspected definitions and hypotheses; they are not human peer review.\par}

\paragraph{What preservation checks mean.}
The final packet contains source, task history, logs, claims, this report and
historical frozen artifacts. An internal manifest hashes every payload entry;
verification checks the archive bytes and recorded evidence references.
Hash/ZIP validation does not compile Lean, prove mathematics, authenticate
an author, or replace independent reading. Dependency runtimes and private
working caches are excluded. The Chinese project report documents the full
route, failed assumptions, recovery protocol and bounded future work packages.
The delivery guide states exact contents, commands and intentional exclusions.

\paragraph{Closure and optional continuation.}
The user selected research-delivery closure: no further mathematical task is
required for this phase. A later, separately scoped continuation could target
higher standard Tor, actual semigroup reductions, or sharper extremal constants.
Each needs its own precise hypotheses, tests and new version evidence.
No such extension, external contact or publication is initiated by this delivery.
A bounded earlier primary-source search did not locate the same conclusion;
that fact does not certify novelty. Versions 0.1 and 0.2 remain historical
snapshots of their earlier proof boundaries.

\begin{thebibliography}{9}
\bibitem{CMS} G.~Caviglia, A.~Moscariello and A.~Sammartano,
\emph{Bounds for syzygies of monomial curves},
Proc. Amer. Math. Soc. \textbf{152} (2024), 3665--3678.
\href{https://doi.org/10.1090/proc/16862}{doi:10.1090/proc/16862}.
Version used: \href{https://arxiv.org/abs/2307.05770v2}{arXiv:2307.05770v2}, 20 July 2024.
\bibitem{HS} J.~Herzog and D.~I.~Stamate,
\emph{On the defining equations of the tangent cone of a numerical semigroup ring},
J. Algebra \textbf{418} (2014), 8--28.
\href{https://doi.org/10.1016/j.jalgebra.2014.07.008}{doi:10.1016/j.jalgebra.2014.07.008}.
Version used: \href{https://arxiv.org/abs/1308.4644v3}{arXiv:1308.4644v3}, 29 July 2014.
\end{thebibliography}
\end{document}
